\chapter{2014年浙江卷（文）}

\section{单选题}

\begin{problem}
设集合 \( S = \{x \mid x \geqslant 2\} \)，\( T = \{x \mid x \leqslant 5\} \)，则 \( S \cap T = \)（\quad）

\choices
    {\((- \infty, 5]\)}
    {\([2, +\infty)\)}
    {\((2,5)\)}
    {\([2,5]\)}
\end{problem}

\begin{answer}
D
\end{answer}

\begin{solution}
由集合的定义，\(S\) 中的元素满足 \(x\geqslant2\)，\(T\) 中的元素满足 \(x\leqslant5\). 因而同时满足这两个条件的 \(x\) 满足

\[
2\leqslant x\leqslant5
\]

所以 \(S\cap T=[2,5]\)，选择 D.
\end{solution}

\begin{problem}
设四边形 \(ABCD\) 的两条对角线为 \(AC\)，\(BD\)，则“四边形 \(ABCD\) 为菱形”是 “\(AC \perp BD\)”的（\quad）

\choices
    {充分不必要条件}
    {必要不充分条件}
    {充分必要条件}
    {既不充分也不必要条件}
\end{problem}

\begin{answer}
A
\end{answer}

\begin{solution}
若四边形 \(ABCD\) 是菱形，则其两条对角线互相垂直，因此“四边形 \(ABCD\) 为菱形”能推出“\(AC\perp BD\)”.

反之，两条对角线垂直的四边形不一定是菱形. 例如取
\(A(-2,0)\)，\(B(0,1)\)，\(C(2,0)\)，\(D(0,-2)\)，则 \(AC\) 为水平线段，\(BD\) 为竖直线段，故 \(AC\perp BD\). 又 \(AB=BC=\sqrt{5}\)，\(CD=DA=\sqrt{8}\)，所以四边形 \(ABCD\) 不是菱形. 因此该条件是充分不必要条件，选择 A.
\end{solution}

\begin{problem}
某几何体的三视图（单位：\(\mathrm{cm}\)）如图所示，则该几何体的体积是（\quad）

\bitmapfigure[max width=0.62\linewidth]{img/2014/zhejiang_liberal/q3_fig1.png}

\choices
    {\(72 \mathrm{~cm}^{3}\)}
    {\(90 \mathrm{~cm}^{3}\)}
    {\(108 \mathrm{~cm}^{3}\)}
    {\(138 \mathrm{~cm}^{3}\)}
\end{problem}

\begin{answer}
B
\end{answer}

\begin{solution}
由三视图可知，该几何体可以分成一个直三棱柱和一个长方体.

直三棱柱的底面是直角三角形，两直角边长分别为 \(4\mathrm{~cm}\) 和 \(3\mathrm{~cm}\)，棱柱的高为 \(3\mathrm{~cm}\)，故其体积为

\[
V_{1}=\frac{1}{2}\times4\times3\times3=18\mathrm{~cm}^{3}
\]

长方体的长、宽、高分别为 \(4\mathrm{~cm}\)，\(6\mathrm{~cm}\)，\(3\mathrm{~cm}\)，故其体积为

\[
V_{2}=4\times6\times3=72\mathrm{~cm}^{3}
\]

这里正视图中的三角形部分沿侧视图中标出的 \(3\mathrm{~cm}\) 延伸，长方体则沿侧视图的总长度 \(6\mathrm{~cm}\) 延伸. 因而

\[
V=V_{1}+V_{2}=18+72=90\mathrm{~cm}^{3}
\]

所以选择 B.
\end{solution}

\begin{problem}
为了得到函数 \( y = \sin 3x + \cos 3x \) 的图象，可以将函数 \( y = \sqrt{2}\cos 3x \) 的图象（\quad）

\choices
    {向右平移 \(\frac{\pi}{4}\) 个单位}
    {向左平移 \(\frac{\pi}{4}\) 个单位}
    {向右平移 \(\frac{\pi}{12}\) 个单位}
    {向左平移 \(\frac{\pi}{12}\) 个单位}
\end{problem}

\begin{answer}
C
\end{answer}

\begin{solution}
利用和角公式，得

\[
\sin3x+\cos3x
=\sqrt{2}\left(\frac{1}{\sqrt{2}}\sin3x+\frac{1}{\sqrt{2}}\cos3x\right)
=\sqrt{2}\cos\left(3x-\frac{\pi}{4}\right)
\]

又

\[
3x-\frac{\pi}{4}=3\left(x-\frac{\pi}{12}\right)
\]

所以目标函数的图象为将 \(y=\sqrt{2}\cos3x\) 的图象向右平移 \(\frac{\pi}{12}\) 个单位所得. 选择 C.
\end{solution}

\begin{problem}
已知圆 \(x^{2} + y^{2} + 2x - 2y + a = 0\) 截直线 \(x + y + 2 = 0\) 所得弦的长度为 \(4\)，则实数 \(a\) 的值为（\quad）

\choices
    {\(-2\)}
    {\(-4\)}
    {\(-6\)}
    {\(-8\)}
\end{problem}

\begin{answer}
B
\end{answer}

\begin{solution}
将圆的方程配方，得

\[
(x+1)^{2}+(y-1)^{2}=2-a
\]

所以圆心为 \(O(-1,1)\)，半径的平方为 \(r^{2}=2-a\). 圆心到直线 \(x+y+2=0\) 的距离为

\[
d=\frac{|-1+1+2|}{\sqrt{1^{2}+1^{2}}}=\sqrt{2}
\]

弦长为 \(4\)，弦的一半长为 \(2\). 由圆心、弦中点和弦端点构成的直角三角形，得

\[
r^{2}=d^{2}+2^{2}=2+4=6
\]

因此 \(2-a=6\)，解得 \(a=-4\)，选择 B.
\end{solution}

\begin{problem}
设 \(m\)，\(n\) 是两条不同的直线，\(\alpha\)，\(\beta\) 是两个不同的平面（\quad）

\choices
    {若 \(m \perp n\)，\(n \parallel \alpha\)，则 \(m \perp \alpha\)}
    {若 \(m \parallel \beta\)，\(\beta \perp \alpha\)，则 \(m \perp \alpha\)}
    {若 \(m \perp \beta\)，\(n \perp \beta\)，\(n \perp \alpha\)，则 \(m \perp \alpha\)}
    {若 \(m \perp n\)，\(n \perp \beta\)，\(\beta \perp \alpha\)，则 \(m \perp \alpha\)}
\end{problem}

\begin{answer}
C
\end{answer}

\begin{solution}
选项 A 不成立. 例如取 \(\alpha:z=0\)，令 \(n\) 为平面 \(z=1\) 内的 \(x\) 方向直线，令 \(m\) 为过 \(n\) 上一点且沿 \(y\) 方向的直线. 则 \(m\perp n\)，\(n\parallel\alpha\)，但 \(m\) 的方向在 \(\alpha\) 内，故 \(m\not\perp\alpha\).

选项 B 不成立. 取 \(\alpha:z=0\)，\(\beta:x=0\)，令 \(m\) 为平面 \(x=1\) 内沿 \(y\) 方向的直线，则 \(m\parallel\beta\)，但 \(m\) 的方向在 \(\alpha\) 内，所以 \(m\not\perp\alpha\).

选项 C 成立. 由 \(m\perp\beta\)，\(n\perp\beta\)，可知 \(m\parallel n\). 又 \(n\perp\alpha\)，所以与 \(n\) 平行的直线 \(m\) 也垂直于 \(\alpha\)，即 \(m\perp\alpha\).

选项 D 不成立. 仍取互相垂直的平面 \(\alpha\) 和 \(\beta\)，令 \(n\perp\beta\)，再在 \(\alpha\) 内取与 \(n\) 垂直的直线 \(m\). 此时 \(m\perp n\)，\(n\perp\beta\)，\(\beta\perp\alpha\)，但 \(m\subset\alpha\)，故 \(m\not\perp\alpha\). 所以选择 C.
\end{solution}

\begin{problem}
已知函数 \( f(x) = x^{3} + ax^{2} + bx + c \)，且 \( 0 < f(-1) = f(-2) = f(-3) \leqslant 3 \)，则（\quad）

\choices
    {\(c \leqslant 3\)}
    {\(3 < c \leqslant 6\)}
    {\(6 < c \leqslant 9\)}
    {\(c > 9\)}
\end{problem}

\begin{answer}
C
\end{answer}

\begin{solution}
由 \(f(-1)=f(-2)\)，得

\[
-1+a-b+c=-8+4a-2b+c
\]

即 \(b=3a-7\). 由 \(f(-2)=f(-3)\)，得

\[
-8+4a-2b+c=-27+9a-3b+c
\]

即 \(b=5a-19\). 联立两式，得 \(a=6\)，\(b=11\).

于是

\[
f(-1)=-1+6-11+c=c-6
\]

由题设 \(0<f(-1)\leqslant3\)，得

\[
0<c-6\leqslant3
\]

即 \(6<c\leqslant9\). 所以选择 C.
\end{solution}

\begin{problem}
在同一直角坐标系中，函数 \( f(x) = x^a (x \geqslant 0) \)，\( g(x) = \log_a x \) 的图象可能是（\quad）

\choices
    {\choicebitmap[width=0.15\paperwidth]{img/2014/zhejiang_liberal/q8_choice_a.png}}
    {\choicebitmap[width=0.15\paperwidth]{img/2014/zhejiang_liberal/q8_choice_b.png}}
    {\choicebitmap[width=0.15\paperwidth]{img/2014/zhejiang_liberal/q8_choice_c.png}}
    {\choicebitmap[width=0.15\paperwidth]{img/2014/zhejiang_liberal/q8_choice_d.png}}
\end{problem}

\begin{answer}
D
\end{answer}

\begin{solution}
因为 \(f(x)=x^{a}\) 在 \(x>0\) 上对任意允许的底数 \(a>0\)、\(a\ne1\) 都是增函数，并且经过原点，所以选项 A 不可能.

若 \(a>1\)，则 \(f(x)=x^{a}\) 在 \(x>0\) 上是增函数且图象向上凸，\(g(x)=\log_{a}x\) 是增函数；若 \(0<a<1\)，则 \(f(x)=x^{a}\) 是增函数且图象向下凹，而 \(g(x)=\log_{a}x\) 是减函数，并经过点 \((1,0)\). 只有选项 D 同时具有“幂函数经过原点且向上增加”和“底数小于 \(1\) 的对数函数递减并经过 \((1,0)\)”的特征. 所以选择 D.
\end{solution}

\begin{problem}
设 \(\theta\) 为两个非零向量 \(\bs{a}\)，\(\bs{b}\) 的夹角. 已知对任意实数 \(t\)，\(|\bs{b} + t\bs{a}|\) 的最小值为 \(1\). 则（\quad）

\choices
    {若 \(\theta\) 确定，则 \(|\bs{a}|\) 唯一确定}
    {若 \(\theta\) 确定，则 \(|\bs{b}|\) 唯一确定}
    {若 \(|\bs{a}|\) 确定，则 \(\theta\) 唯一确定}
    {若 \(|\bs{b}|\) 确定，则 \(\theta\) 唯一确定}
\end{problem}

\begin{answer}
B
\end{answer}

\begin{solution}
设 \(A=|\bs{a}|>0\)，\(B=|\bs{b}|>0\). 则

\[
|\bs{b}+t\bs{a}|^{2}=B^{2}+2tAB\cos\theta+t^{2}A^{2}
\]

这是关于 \(t\) 的二次函数，其最小值在

\[
t=-\frac{B\cos\theta}{A}
\]

处取得，且

\[
\min|\bs{b}+t\bs{a}|^{2}=B^{2}-B^{2}\cos^{2}\theta=B^{2}\sin^{2}\theta
\]

题设最小值为 \(1>0\)，所以 \(\bs{a}\)，\(\bs{b}\) 不共线，从而 \(0<\theta<\pi\). 因此

\[
B\sin\theta=1
\]

因此当 \(\theta\) 确定时，\(B=|\bs{b}|=\frac{1}{\sin\theta}\) 唯一确定，选项 B 正确. 选项 A 不正确：取 \(\theta=\frac{\pi}{2}\)，\(B=1\)，则任意 \(A>0\) 都满足最小值为 \(1\)，故 \(A\) 不能唯一确定. 选项 C 不正确：固定 \(A=1\)，取 \((B,\theta)=(2,\frac{\pi}{6})\) 或 \((2,\frac{5\pi}{6})\)，均有 \(B\sin\theta=1\)，但 \(\theta\) 不同. 选项 D 不正确：固定 \(B=2\)，上述两个不同的 \(\theta\) 都满足条件，故 \(\theta\) 不能唯一确定.
\end{solution}

\begin{problem}
如图，某人在垂直于水平地面 \(ABC\) 的墙面前的点 \(A\) 处进行射击训练，已知点 \(A\) 到墙面的距离为 \(AB\)，某目标点 \(P\) 沿墙面上的射线 \(CM\) 移动，此人为了准确瞄准目标点 \(P\)，需计算由点 \(A\) 观察点 \(P\) 的仰角 \(\theta\) 的大小（仰角 \(\theta\) 为直线 \(AP\) 与平面 \(ABC\) 所成的角）. 若 \(AB = 15 \mathrm{~m}\)，\(AC = 25 \mathrm{~m}\)，\(\angle BCM = 30^{\circ}\)，则 \(\tan \theta\) 的最大值是（\quad）

\bitmapfigure[max width=0.36\linewidth]{img/2014/zhejiang_liberal/q10_fig1.png}

\choices
    {\(\frac{\sqrt{30}}{5}\)}
    {\(\frac{\sqrt{30}}{10}\)}
    {\(\frac{4\sqrt{3}}{9}\)}
    {\(\frac{5\sqrt{3}}{9}\)}
\end{problem}

\begin{answer}
D
\end{answer}

\begin{solution}
设 \(Q\) 为点 \(P\) 在水平地面 \(ABC\) 上的射影，令 \(x=CQ\). 由图中 \(\angle BCM=30^{\circ}\)，目标点沿射线 \(CM\) 移动，故

\[
PQ=x\tan30^{\circ}=\frac{x}{\sqrt{3}}
\]

在水平地面内，\(AB\perp BC\)，且

\[
BC=\sqrt{AC^{2}-AB^{2}}=\sqrt{25^{2}-15^{2}}=20
\]

因而

\[
AQ^{2}=AB^{2}+(20-x)^{2}=225+(20-x)^{2}
\]

直线 \(AP\) 与地面所成的角为 \(\theta\)，所以

\[
\tan^{2}\theta=\frac{PQ^{2}}{AQ^{2}}=\frac{x^{2}}{3\left(225+(20-x)^{2}\right)}
\]

令

\[
h(x)=\frac{x^{2}}{225+(20-x)^{2}}
\]

则 \(x\geqslant0\)，并且

\[
h'(x)=\frac{2x(625-20x)}{\left(225+(20-x)^{2}\right)^{2}}
\]

所以 \(h(x)\) 在 \(x=\frac{125}{4}\) 处取得最大值. 代入得

\[
\tan^{2}\theta_{\max}=\frac{1}{3}\cdot\frac{(125/4)^{2}}{225+(-45/4)^{2}}=\frac{25}{27}
\]

由于 \(\theta\) 为仰角，\(\tan\theta>0\)，故

\[
\tan\theta_{\max}=\frac{5}{\sqrt{27}}=\frac{5\sqrt{3}}{9}
\]

所以选择 D.
\end{solution}

\section{填空题}

\begin{problem}
已知 \(\i\) 是虚数单位，计算 \( \frac{1-\i}{(1+\i)^{2}} =\fillinblank{} \).
\end{problem}

\begin{answer}
\(\frac{-1-\i}{2}\)
\end{answer}

\begin{solution}
先计算分母：

\[
(1+\i)^{2}=1+2\i+\i^{2}=2\i
\]

因此

\[
\frac{1-\i}{(1+\i)^{2}}=\frac{1-\i}{2\i}=\frac{(1-\i)(-\i)}{2}=\frac{-1-\i}{2}
\]
\end{solution}

\begin{problem}
若实数 \(x\)，\(y\) 满足 \(\begin{cases}
x + 2y - 4 \leqslant 0,\\
x - y - 1 \leqslant 0,\\
x \geqslant 1,
\end{cases}\) 则 \(x + y\) 的取值范围是\fillinblank{}.
\end{problem}

\begin{answer}
\([1,3]\)
\end{answer}

\begin{solution}
由约束条件，得

\(y\leqslant2-\frac{x}{2}\)，\(y\geqslant x-1\)，\(x\geqslant1\).

要使 \(y\) 的上下界有交集，需满足

\[
x-1\leqslant2-\frac{x}{2}
\]

即 \(1\leqslant x\leqslant2\). 对固定的 \(x\)，有

\[
2x-1\leqslant x+y\leqslant2+\frac{x}{2}
\]

当 \(x=1\) 且 \(y=0\) 时，\(x+y=1\)；当 \(x=2\) 且 \(y=1\) 时，\(x+y=3\). 又线段

\((x,y)=(1,0)+t(1,1)\)，\(0\leqslant t\leqslant1\)

上的点均满足约束条件，且 \(x+y=1+2t\) 能取遍 \([1,3]\). 所以 \(x+y\) 的取值范围是 \([1,3]\).
\end{solution}

\begin{problem}
若某程序框图如图所示，当输入 \(50\) 时，则该程序运行后输出的结果是\fillinblank{}.

\texfigure[float=inline,max width=0.38\linewidth]{img_repaint/2014/zhejiang_liberal/q13_fig1.tex}
\end{problem}

\begin{answer}
\(6\)
\end{answer}

\begin{solution}
按照程序框图逐次计算. 初始 \(S=0\)，\(i=1\)，每次先执行 \(S=2S+i\)，再执行 \(i=i+1\). 第一次计算后 \(S=1\)，\(i=2\)；第二次计算后 \(S=4\)，\(i=3\)；第三次计算后 \(S=11\)，\(i=4\)；第四次计算后 \(S=26\)，\(i=5\)；第五次计算后 \(S=57\)，\(i=6\).

前四次计算后 \(S\leqslant50\)，第五次计算后 \(S=57>50\)，此时输出更新后的 \(i=6\). 所以输出结果为 \(6\).
\end{solution}

\begin{problem}
在 \(3\) 张奖券中有一、二等奖各一张，另 \(1\) 张无奖. 甲，乙两人各抽取一张，两人都中奖的概率为\fillinblank{}.
\end{problem}

\begin{answer}
\(\frac{1}{3}\)
\end{answer}

\begin{solution}
将一等奖、二等奖和无奖分别记为三张不同的奖券. 甲先抽、乙后抽时，所有等可能的有序结果共有

\[
3\times2=6
\]

种. 两人都中奖时，甲、乙抽到一等奖和二等奖，或甲、乙抽到二等奖和一等奖，共 \(2\) 种. 因而所求概率为

\[
\frac{2}{6}=\frac{1}{3}
\]
\end{solution}

\begin{problem}
设函数 \( f(x) = \begin{cases}
x^2 + 2x + 2, & x \leqslant 0,\\
-x^2, & x > 0.
\end{cases} \) 若 \( f(f(a)) = 2 \)，则 \( a = \)\fillinblank{}.
\end{problem}

\begin{answer}
\(\sqrt{2}\)
\end{answer}

\begin{solution}
令 \(u=f(a)\). 要使 \(f(u)=2\)，若 \(u>0\)，则 \(f(u)=-u^{2}<0\)，不可能等于 \(2\). 因此必须有 \(u\leqslant0\)，从而

\[
u^{2}+2u+2=2
\]

即

\(u(u+2)=0\)，\(u=0\text{ 或 }u=-2\).

下面分别考察 \(f(a)=0\) 和 \(f(a)=-2\). 当 \(a\leqslant0\) 时，\(f(a)=(a+1)^{2}+1\geqslant1\)，所以这两种情况都不可能由 \(a\leqslant0\) 得到；当 \(a>0\) 时，\(f(a)=-a^{2}\). 方程 \(-a^{2}=0\) 没有满足 \(a>0\) 的解，而

\(-a^{2}=-2\)，\(\Longrightarrow\)，\(a^{2}=2\)，\(\Longrightarrow\)，\(a=\sqrt{2}\).

所以 \(a=\sqrt{2}\).
\end{solution}

\begin{problem}
已知实数 \(a\)，\(b\)，\(c\) 满足 \(a + b + c = 0\)，\(a^2 + b^2 + c^2 = 1\)，则 \(a\) 的最大值是\fillinblank{}.
\end{problem}

\begin{answer}
\(\frac{\sqrt{6}}{3}\)
\end{answer}

\begin{solution}
由 \(a+b+c=0\)，得 \(b+c=-a\). 由平方和不小于平均值关系，

\[
b^{2}+c^{2}\geqslant\frac{(b+c)^{2}}{2}=\frac{a^{2}}{2}
\]

因此

\[
1=a^{2}+b^{2}+c^{2}\geqslant a^{2}+\frac{a^{2}}{2}=\frac{3a^{2}}{2}
\]

从而 \(a^{2}\leqslant\frac{2}{3}\)，所以

\[
a\leqslant\sqrt{\frac{2}{3}}=\frac{\sqrt{6}}{3}
\]

当 \(b=c=-\frac{a}{2}\)，且 \(a=\sqrt{\frac{2}{3}}\) 时，两个条件均成立，故上界可以取得. 所以 \(a\) 的最大值为 \(\frac{\sqrt{6}}{3}\).
\end{solution}

\begin{problem}
设直线 \(x - 3y + m = 0\)（\(m \neq 0\)）与双曲线 \(\frac{x^2}{a^2} - \frac{y^2}{b^2} = 1\)（\(a > 0\)，\(b > 0\)）的两条渐近线分别交于点 \(A\)，\(B\). 若点 \(P(m, 0)\) 满足 \(|PA| = |PB|\)，则该双曲线的离心率是\fillinblank{}.
\end{problem}

\begin{answer}
\(\frac{\sqrt{5}}{2}\)
\end{answer}

\begin{solution}
设双曲线渐近线的斜率为 \(k=\frac{b}{a}>0\)，则两条渐近线为 \(y=kx\) 和 \(y=-kx\). 由于给定直线分别与两条渐近线相交，\(3k-1\ne0\). 由联立方程可得交点分别为 \(A\left(\frac{m}{3k-1},\frac{km}{3k-1}\right)\) 和 \(B\left(-\frac{m}{3k+1},\frac{km}{3k+1}\right)\).

由于 \(P=(m,0)\)，

于是 \(\frac{|PA|^{2}}{m^{2}}=\frac{(3k-2)^{2}+k^{2}}{(3k-1)^{2}}\)，且 \(\frac{|PB|^{2}}{m^{2}}=\frac{(3k+2)^{2}+k^{2}}{(3k+1)^{2}}\).

由 \(|PA|=|PB|\)，并约去非零因子 \(m^{2}\)，得

\[
\frac{5k^{2}-6k+2}{(3k-1)^{2}}=\frac{5k^{2}+6k+2}{(3k+1)^{2}}
\]

交叉相乘并化简，得

\[
12k(1-4k^{2})=0
\]

因为 \(k>0\)，所以 \(k=\frac{1}{2}\)，即 \(\frac{b}{a}=\frac{1}{2}\). 双曲线的离心率满足

\[
e=\sqrt{1+\frac{b^{2}}{a^{2}}}=\sqrt{1+k^{2}}=\sqrt{1+\frac14}=\frac{\sqrt{5}}{2}
\]
\end{solution}

\section{解答题}

\begin{problem}
在 \(\triangle ABC\) 中，内角 \(A\)，\(B\)，\(C\) 所对的边分别为 \(a\)，\(b\)，\(c\). 已知 \(4\sin^2\frac{A - B}{2} + 4\sin A\sin B = 2 + \sqrt{2}\).

\begin{enumerate}
\item 求角 \(C\) 的大小；
\item 已知 \(b = 4\)，\(\triangle ABC\) 的面积为 6，求边长 \(c\) 的值.
\end{enumerate}
\end{problem}

\begin{solution}
\begin{enumerate}
\item 由 \(A+B+C=\pi\)，以及倍角公式，有

\[
4\sin^2\frac{A-B}{2}=2-2\cos(A-B)
\]

\[
4\sin A\sin B=2\cos(A-B)-2\cos(A+B)
\]

因此，已知等式左边为

\[
2-2\cos(A+B)=2-2\cos(\pi-C)=2+2\cos C
\]

所以 \(2+2\cos C=2+\sqrt{2}\)，即 \(\cos C=\frac{\sqrt{2}}{2}\). 因为 \(0<C<\pi\)，故 \(C=\frac{\pi}{4}\).

\item 由三角形面积公式，得

\[
6=\frac{1}{2}ab\sin C=\frac{1}{2}\cdot a\cdot4\cdot\frac{\sqrt{2}}{2}=a\sqrt{2}
\]

所以 \(a=3\sqrt{2}\). 再由余弦定理，得

\[
c^2=a^2+b^2-2ab\cos C=18+16-2\cdot3\sqrt{2}\cdot4\cdot\frac{\sqrt{2}}{2}=10
\]

由于 \(c>0\)，所以 \(c=\sqrt{10}\).
\end{enumerate}
\end{solution}

\begin{problem}
已知等差数列 \(\{a_{n}\}\) 的公差 \(d > 0\)，设 \(\{a_{n}\}\) 的前 \(n\) 项和为 \(S_{n}\)，\(a_{1} = 1\)，\(S_{2} \cdot S_{3} = 36\).

\begin{enumerate}
\item 求 \(d\) 及 \(S_{n}\)；
\item 求 \(m\)，\(k\)（\(m\)，\(k \in \N^*\)）的值，使得 \(a_m + a_{m+1} + a_{m+2} + \cdots + a_{m+k} = 65\).
\end{enumerate}
\end{problem}

\begin{solution}
\begin{enumerate}
\item 由等差数列求和公式，得 \(S_2=2+d\)，\(S_3=3+3d\).

于是

\[
(d+2)(3d+3)=36
\]

即 \((d+2)(d+1)=12\)，从而 \(d^2+3d-10=0\). 解得 \(d=2\) 或 \(d=-5\). 因为 \(d>0\)，所以 \(d=2\). 因此

\[
a_n=a_1+(n-1)d=2n-1
\]

\[
S_n=\frac{n}{2}\bigl(2a_1+(n-1)d\bigr)=\frac{n}{2}\bigl(2+2n-2\bigr)=n^2
\]

\item 由等差数列求和公式，得

\[
a_m+a_{m+1}+\cdots+a_{m+k}=\frac{k+1}{2}(a_m+a_{m+k})=(k+1)(2m+k-1)
\]

所以

\[
(k+1)(2m+k-1)=65
\]

令 \(u=k+1\)，则 \(u\geqslant2\)，且

\[
u(2m+u-2)=65
\]

由于 \(65=5\cdot13\)，并且 \(u\geqslant2\)，逐一考察正因数可得：当 \(u=5\) 时，\(2m+3=13\)，解得 \(m=5\)；当 \(u=13\) 或 \(u=65\) 时，分别得到 \(m=-3\) 或 \(m=-31\)，均不符合 \(m\in\N^*\). 因而 \(u=5\)，即 \(m=5\)，\(k=4\).
\end{enumerate}
\end{solution}

\begin{problem}
如图，在四棱锥 \(A - BCDE\) 中，平面 \(ABC \perp\) 平面 \(BCDE\)，\(\angle CDE = \angle BED = 90^\circ\)，\(AB = CD = 2\)，\(DE = BE = 1\)，\(AC = \sqrt{2}\).

\begin{enumerate}
\item 证明：\(AC \perp\) 平面 \(BCDE\)；
\item 求直线 \(AE\) 与平面 \(ABC\) 所成的角的正切值.
\end{enumerate}
\bitmapfigure[max width=0.42\linewidth]{img/2014/zhejiang_liberal/q20_fig1.png}
\end{problem}

\begin{solution}
\begin{enumerate}
\item 在平面 \(BCDE\) 内建立直角坐标系，取 \(D(0,0)\)，\(C(2,0)\)，\(E(0,1)\)，\(B(1,1)\).

这与 \(CD=2\)，\(DE=BE=1\) 以及两个直角条件相符，所以 \(BC=\sqrt{2}\). 因而

\[
AB^2=4=2+2=AC^2+BC^2
\]

由勾股定理的逆定理，\(\angle ACB=90^\circ\)，即 \(AC\perp BC\). 又平面 \(ABC\) 与平面 \(BCDE\) 的交线为 \(BC\)，所以由二面角的性质，得 \(AC\perp\) 平面 \(BCDE\).

\item 在上述坐标系中，直线 \(BC\) 的方程为 \(x+y=2\)，直线 \(BD\) 的方程为 \(y=x\)，故 \(BD\perp BC\). 过点 \(E\) 作 \(EF\parallel BD\)，交直线 \(BC\) 于 \(F\). 于是 \(EF\) 的方程为 \(y=x+1\)，从而

\[
F\left(\frac{1}{2},\frac{3}{2}\right)
\]

所以 \(EF=BF=\frac{1}{\sqrt{2}}\)，\(CF=\frac{3}{\sqrt{2}}\).

由 \(BD\perp BC\) 且 \(BD\) 在平面 \(BCDE\) 内，结合两平面垂直，得 \(BD\perp\) 平面 \(ABC\)，从而 \(EF\perp\) 平面 \(ABC\). 因为 \(F\in BC\)，所以 \(AF\) 是 \(AE\) 在平面 \(ABC\) 内的射影，\(\angle EAF\) 就是直线 \(AE\) 与平面 \(ABC\) 所成的角.

又因为 \(AC\perp\) 平面 \(BCDE\)，所以 \(AC\perp CF\)，在直角三角形 \(ACF\) 中

\[
AF^2=AC^2+CF^2=2+\frac{9}{2}=\frac{13}{2}
\]

因此

\[
\tan\angle EAF=\frac{EF}{AF}=\frac{1/\sqrt{2}}{\sqrt{13/2}}=\frac{1}{\sqrt{13}}
\]
\end{enumerate}
\end{solution}

\begin{problem}
已知函数 \( f(x) = x^3 + 3|x - a| \) （\( a > 0 \)），若 \( f(x) \) 在 \([-1, 1]\) 上的最小值记为 \( g(a) \).

\begin{enumerate}
\item 求 \(g(a)\)；
\item 证明：当 \(x \in [-1, 1]\) 时，恒有 \(f(x) \leqslant g(a) + 4\).
\end{enumerate}
\end{problem}

\begin{solution}
\begin{enumerate}
\item 当 \(0<a\leqslant1\) 时，在区间 \([-1,a]\) 上

\[
f(x)=x^3-3x+3a
\]

其导数为 \(3x^2-3\leqslant0\)，所以该函数在 \([-1,a]\) 上单调递减；在区间 \([a,1]\) 上

\[
f(x)=x^3+3x-3a
\]

其导数为 \(3x^2+3>0\)，所以该函数在 \([a,1]\) 上单调递增. 因而最小值在 \(x=a\) 处取得，\(g(a)=a^3\).

当 \(a>1\) 时，对任意 \(x\in[-1,1]\) 都有 \(x<a\)，故

\[
f(x)=x^3-3x+3a
\]

其导数在 \([-1,1]\) 上不大于零，所以最小值在 \(x=1\) 处取得，\(g(a)=1-3+3a=3a-2\). 综上，

\[
g(a)=\begin{cases}
a^3,&0<a\leqslant1,\\
3a-2,&a>1.
\end{cases}
\]

\item 当 \(a>1\) 时，\(f(x)=x^3-3x+3a\) 在 \([-1,1]\) 上单调递减，所以

\[
f(x)\leqslant f(-1)=3a+2=g(a)+4
\]

当 \(0<a\leqslant1\) 时，若 \(x\in[-1,a]\)，由前面的单调性知

\[
f(x)\leqslant f(-1)=3a+2
\]

又因为

\[
a^3+4-(3a+2)=a^3-3a+2=(a-1)^2(a+2)\geqslant0
\]

所以 \(f(x)\leqslant a^3+4=g(a)+4\). 若 \(x\in[a,1]\)，则

\[
f(x)\leqslant f(1)=4-3a\leqslant a^3+4=g(a)+4
\]

其中最后一个不等式由 \(a>0\) 得出. 因此对任意 \(x\in[-1,1]\)，恒有 \(f(x)\leqslant g(a)+4\).
\end{enumerate}
\end{solution}

\begin{problem}
已知 \(\triangle ABP\) 的三个顶点都在抛物线 \(C: x^2 = 4y\) 上，\(F\) 为抛物线 \(C\) 的焦点，点 \(M\) 为 \(AB\) 的中点，\(\overrightarrow{PF} = 3\overrightarrow{FM}\).

\begin{enumerate}
\item 若 \(|PF| = 3\)，求点 \(M\) 的坐标；
\item 求 \(\triangle ABP\) 面积的最大值.
\end{enumerate}
\bitmapfigure[max width=0.42\linewidth]{img/2014/zhejiang_liberal/q22_fig1.png}
\end{problem}

\begin{solution}
\begin{enumerate}
\item 由抛物线方程 \(x^2=4y\)，可知其焦点为 \(F(0,1)\). 设 \(P\left(t,\frac{t^2}{4}\right)\)，则

\[
PF^2=t^2+\left(\frac{t^2}{4}-1\right)^2=\left(\frac{t^2}{4}+1\right)^2
\]

因为 \(\frac{t^2}{4}+1>0\)，所以 \(PF=\frac{t^2}{4}+1\). 又由 \(PF=3\)，得 \(t^2=8\)，即 \(t=\pm2\sqrt{2}\).

由向量关系

\[
\overrightarrow{PF}=F-P=3(M-F)=3\overrightarrow{FM}
\]

得

\[
M=\frac{4F-P}{3}=\left(-\frac{t}{3},\frac{4}{3}-\frac{t^2}{12}\right)
\]

代入 \(t^2=8\)，可得点 \(M\) 的坐标为

\[
M\left(\pm\frac{2\sqrt{2}}{3},\frac{2}{3}\right)
\]

\item 设 \(A\left(u,\frac{u^2}{4}\right)\)，\(B\left(v,\frac{v^2}{4}\right)\)，仍设 \(P\left(t,\frac{t^2}{4}\right)\). 因为 \(M\) 是 \(AB\) 的中点，且由已知向量关系

\[
M=\left(-\frac{t}{3},\frac{4}{3}-\frac{t^2}{12}\right)
\]

所以 \(u+v=-\frac{2t}{3}\)，\(u^2+v^2=\frac{32}{3}-\frac{2t^2}{3}\).

于是

\[
uv=\frac{(u+v)^2-(u^2+v^2)}{2}=\frac{5t^2}{9}-\frac{16}{3}
\]

从而

\[
(u-v)^2=(u+v)^2-4uv=\frac{16}{9}(12-t^2)
\]

因为 \(u\)，\(v\) 为实数，必须有 \(t^2\leqslant12\). 另一方面，由三点在抛物线 \(y=\frac{x^2}{4}\) 上，三角形面积为

\[
S_{\triangle ABP}=\frac18\left|(u-v)(v-t)(t-u)\right|
\]

又

\[
(u-t)(v-t)=uv-t(u+v)+t^2=\frac{4}{9}(5t^2-12)
\]

因此

\[
S_{\triangle ABP}=\frac{2}{27}\left|5t^2-12\right|\sqrt{12-t^2}
\]

令 \(z=t^2\)，则 \(0\leqslant z\leqslant12\)，只需研究

\[
h(z)=\left|5z-12\right|\sqrt{12-z}
\]

当 \(0\leqslant z\leqslant\frac{12}{5}\) 时，

\[
h(z)=(12-5z)\sqrt{12-z}
\]
其导数为 \(h'(z)=\frac{15z-132}{2\sqrt{12-z}}<0\).

所以这一区间的最大值为 \(h(0)=24\sqrt{3}\). 当 \(\frac{12}{5}\leqslant z<12\) 时，

\[
h(z)=(5z-12)\sqrt{12-z}
\]
其导数为 \(h'(z)=\frac{132-15z}{2\sqrt{12-z}}\).

故 \(h(z)\) 在 \(z=\frac{44}{5}\) 处取得最大值，且

\[
h\left(\frac{44}{5}\right)=32\sqrt{\frac{16}{5}}=\frac{128}{\sqrt{5}}
\]

又 \(\frac{128}{\sqrt{5}}>24\sqrt{3}\)，所以

\[
\max S_{\triangle ABP}=\frac{2}{27}\cdot\frac{128}{\sqrt{5}}=\frac{256}{27\sqrt{5}}=\frac{256\sqrt{5}}{135}
\]

在 \(z=\frac{44}{5}<12\) 时，\((u-v)^2>0\)，故对应的 \(A\)，\(B\) 为不同点，最大值确实可以取得.
\end{enumerate}
\end{solution}
